% =====================================================================
% CAP. 11 - NIVEL 4 - FONTES DEPENDENTES, PROJETO E ESTABILIDADE
% =====================================================================
\blocodesafio

\begin{questao}[4]{}
O circuito RC natural abaixo contém uma VCCS orientada para cima, de valor $g v_C$.
O capacitor está inicialmente carregado. Determine o valor crítico de $g$ que
separa decaimento, resposta constante e crescimento exponencial. Em seguida,
classifique os casos $g=\SI{0.4}{\milli\ensuremath{\mathrm{S}}}$,
$\SI{0.5}{\milli\ensuremath{\mathrm{S}}}$ e $\SI{0.6}{\milli\ensuremath{\mathrm{S}}}$ para
$R=\SI{2}{\kilo\ohm}$ e $C=\SI{100}{\micro\farad}$.

\circuitoqq{%
\begin{circuitikz}[scale=0.98,transform shape,line width=0.8pt,every node/.append style={font=\scriptsize}]
\draw (0,0) -- (7.0,0);
\draw (0,3) -- (7.0,3);
\draw (1.8,3) to[R] (1.8,0);
\node[anchor=east] at (1.55,1.5) {$R=\SI{2}{\kilo\ohm}$};
\draw (4.2,3) to[C,v^>={$v_C$}] (4.2,0);
\node[anchor=east] at (3.95,1.5) {$C=\SI{100}{\micro\farad}$};
\draw (7.0,0) to[american controlled current source] (7.0,3);
\node[anchor=west] at (7.25,1.5) {$g\,v_C$};
\end{circuitikz}}{RC com VCCS capaz de reduzir, anular ou inverter o amortecimento.}

\resposta{$g_{\mathrm{crit}}=1/R=\SI{0.5}{\milli\ensuremath{\mathrm{S}}}$; para
\SI{0.4}{\milli\ensuremath{\mathrm{S}}}: estável, $\tau=\SI{1}{\second}$; para
\SI{0.5}{\milli\ensuremath{\mathrm{S}}}: polo em zero; para \SI{0.6}{\milli\ensuremath{\mathrm{S}}}:
instável, crescimento com constante de tempo de magnitude \SI{1}{\second}}
\resolucao{Pela KCL,
\[
C\dot v_C+\frac{v_C}{R}-gv_C=0
\quad\Rightarrow\quad
\dot v_C+\frac{1/R-g}{C}v_C=0.
\]
O polo é $s=-(1/R-g)/C$. O limite ocorre para $g=1/R$, isto é,
$g_{\mathrm{crit}}=\SI{0.5}{\milli\ensuremath{\mathrm{S}}}$. Para $g=0{,}4\,\mathrm{mS}$,
$1/R-g=0{,}1\,\mathrm{mS}$ e
$\tau=C/(1/R-g)=\SI{1}{\second}$: a resposta decai. Para
$g=0{,}5\,\mathrm{mS}$, o polo é zero e, no modelo ideal, a tensão permanece
constante. Para $g=0{,}6\,\mathrm{mS}$, o coeficiente torna-se negativo e o polo é
$+\SI{1}{\per\second}$, produzindo crescimento exponencial.}
\end{questao}

\begin{questao}[4]{}
Deseja-se usar uma VCCS $g v$ para ajustar o amortecimento do RLC em paralelo
abaixo. Determine o valor de $g$ para obter amortecimento crítico. Depois indique
se a resposta será subamortecida ou superamortecida para valores de $g$ menores ou
maiores que o valor calculado, supondo $g\ge0$.

\circuitoqq{%
\begin{circuitikz}[scale=0.94,transform shape,line width=0.8pt,every node/.append style={font=\scriptsize}]
\draw (0,0) -- (8.8,0);
\draw (0,3) -- (8.8,3);
\draw (1.5,3) to[R] (1.5,0);
\node[anchor=east] at (1.25,1.5) {$R=\SI{100}{\ohm}$};
\draw (3.9,3) to[C,v^>={$v$}] (3.9,0);
\node[anchor=east] at (3.65,1.5) {$C=\SI{100}{\micro\farad}$};
\draw (6.3,3) to[L] (6.3,0);
\node[anchor=east] at (6.05,1.5) {$L=\SI{100}{\milli\henry}$};
\draw (8.8,3) to[american controlled current source] (8.8,0);
\node[anchor=west] at (9.05,1.5) {$g\,v$};
\end{circuitikz}}{Projeto do amortecimento de um circuito RLC em paralelo usando uma VCCS.}

\resposta{$g_{\mathrm{crit}}\approx\SI{53.25}{\milli\ensuremath{\mathrm{S}}}$; abaixo desse
valor a resposta é subamortecida; acima, superamortecida}
\resolucao{A equação natural é
\[
C\ddot v+\left(\frac1R+g\right)\dot v+\frac1L v=0.
\]
Logo
\[
\alpha=\frac{1/R+g}{2C},\qquad
\omega_0=\frac1{\sqrt{LC}}.
\]
Com $L=0{,}1\,\mathrm{H}$ e $C=100\,\mathrm{\mu F}$,
$\omega_0\approx\SI{316.23}{\radian\per\second}$. Para amortecimento crítico,
$\alpha=\omega_0$, então
\[
\frac1R+g=2C\omega_0
=2(100\times10^{-6})(316{,}23)\approx\SI{63.25}{\milli\ensuremath{\mathrm{S}}}.
\]
Como $1/R=\SI{10}{\milli\ensuremath{\mathrm{S}}}$,
$g_{\mathrm{crit}}\approx\SI{53.25}{\milli\ensuremath{\mathrm{S}}}$. Reduzir $g$ reduz
$\alpha$ e leva ao caso subamortecido; aumentar $g$ produz superamortecimento.}
\end{questao}

\begin{questao}[4]{}
No RLC em série abaixo, a VCVS tem valor $\mu v_C$ e polaridade tal que sua tensão
\textbf{auxilia} a tensão do capacitor na equação de malha, resultando no termo
restaurador $(1-\mu)v_C$. Para
$R=\SI{20}{\ohm}$, $L=\SI{0.1}{\henry}$ e
$C=\SI{100}{\micro\farad}$, derive a equação característica e determine os
intervalos de $\mu$ que produzem resposta subamortecida, criticamente amortecida,
superamortecida estável, polo em zero e instabilidade.

\circuitoqq{%
\begin{circuitikz}[scale=0.88,transform shape,line width=0.8pt,every node/.append style={font=\scriptsize}]
\draw (0,0) to[R] (2.2,0)
      to[L,i>^={$i$}] (4.8,0)
      to[C,v^>={$v_C$}] (7.4,0)
      to[american controlled voltage source] (10.0,0)
      -- (10.0,-2.2) -- (0,-2.2) -- (0,0);
\node[anchor=south] at (1.1,0.25) {$R=\SI{20}{\ohm}$};
\node[anchor=south] at (3.5,0.25) {$L=\SI{0.1}{\henry}$};
\node[anchor=south] at (6.1,0.25) {$C=\SI{100}{\micro\farad}$};
\node[anchor=south] at (8.7,0.25) {$\mu\,v_C$};
\end{circuitikz}}{RLC com VCVS modificando o termo restaurador e a estabilidade.}

\resposta{$s^2+200s+100000(1-\mu)=0$; subamortecida para $\mu<0{,}9$;
crítica em $\mu=0{,}9$; superamortecida estável para $0{,}9<\mu<1$;
polo em zero para $\mu=1$; instável para $\mu>1$}
\resolucao{Com a polaridade indicada, a KVL é
$L\dot i+Ri+(1-\mu)v_C=0$. Como $i=C\dot v_C$,
\[
LC\ddot v_C+RC\dot v_C+(1-\mu)v_C=0.
\]
Dividindo por $LC$:
\[
s^2+\frac RL s+\frac{1-\mu}{LC}=0
=s^2+200s+100000(1-\mu)=0.
\]
A fronteira entre polos complexos e reais ocorre quando o discriminante é zero:
$200^2-4\cdot100000(1-\mu)=0$, fornecendo $\mu=0{,}9$. Para
$\mu<0{,}9$ há polos complexos conjugados com parte real negativa. Em
$\mu=0{,}9$ há polo real duplo. Para $0{,}9<\mu<1$, os dois polos são reais e
negativos. Em $\mu=1$, o termo constante se anula e surge um polo em zero. Para
$\mu>1$, o produto dos polos é negativo, portanto um polo é positivo e o circuito
é instável.}
\end{questao}

\begin{questao}[4]{}
A rede abaixo possui dois capacitores e uma CCCS que injeta no nó $v_2$ a corrente
$\beta i_x$, em que $i_x=v_1/R_1$. Considere
$R_1=R_2=\SI{1}{\kilo\ohm}$ e $C_1=C_2=\SI{1}{\milli\farad}$.

\circuitoqq{%
\begin{circuitikz}[scale=0.92,transform shape,line width=0.8pt,every node/.append style={font=\scriptsize}]
\draw (0,0) -- (8.4,0);
\draw (0,3) coordinate(a) to[R] (5.0,3) coordinate(b);
\node[anchor=south] at (2.5,3.25) {$R_2=\SI{1}{\kilo\ohm}$};
\draw (a) to[R,i>^={$i_x$}] (0,0);
\node[anchor=west] at (0.25,1.5) {$R_1=\SI{1}{\kilo\ohm}$};
\draw (2.5,3) to[C,v^>={$v_1$}] (2.5,0);
\node[anchor=east] at (2.25,1.5) {$C_1=\SI{1}{\milli\farad}$};
\draw (b) to[C,v^>={$v_2$}] (5.0,0);
\node[anchor=east] at (4.75,1.5) {$C_2=\SI{1}{\milli\farad}$};
\draw (8.4,0) to[american controlled current source] (8.4,3) -- (b);
\node[anchor=west] at (8.65,1.5) {$\beta\,i_x$};
\end{circuitikz}}{Rede RC de segunda ordem com CCCS e possibilidade de perda de estabilidade.}

\begin{itensq}
\item Obtenha as equações de estado em $v_1$ e $v_2$.
\item Determine a equação característica em função de $\beta$.
\item Encontre a faixa de $\beta$ que garante estabilidade assintótica.
\item Para $\beta=0{,}5$, $v_1(0)=\SI{1}{\volt}$ e $v_2(0)=0$, determine $v_2(t)$.
\end{itensq}

\resposta{$\dot v_1=-2v_1+v_2$;
$\dot v_2=(1+\beta)v_1-v_2$;
$s^2+3s+(1-\beta)=0$; estável para $\beta<1$;
para $\beta=0{,}5$,
$v_2(t)=\frac{1.5}{\sqrt7}\left(e^{(-3+\sqrt7)t/2}-e^{(-3-\sqrt7)t/2}\right)\ \si{\volt}$}
\resolucao{Como $R_1C_1=R_2C_1=R_2C_2=\SI{1}{\second}$, a KCL no nó 1 fornece
\[
C_1\dot v_1+\frac{v_1}{R_1}+\frac{v_1-v_2}{R_2}=0
\quad\Rightarrow\quad
\dot v_1=-2v_1+v_2.
\]
No nó 2,
\[
C_2\dot v_2+\frac{v_2-v_1}{R_2}-\beta\frac{v_1}{R_1}=0,
\]
portanto $\dot v_2=(1+\beta)v_1-v_2$. A matriz de estado é
\[
A=\begin{bmatrix}-2&1\\1+\beta&-1\end{bmatrix},
\]
e
$\det(sI-A)=s^2+3s+(1-\beta)$. Como o coeficiente de $s$ é positivo, a
estabilidade assintótica de segunda ordem exige $1-\beta>0$, isto é,
$\beta<1$. Para $\beta=0{,}5$, os polos são
$s_{1,2}=(-3\pm\sqrt7)/2$. Além disso, $v_2(0)=0$ e
$\dot v_2(0)=1{,}5\,\si{\volt\per\second}$. Logo
$v_2=A_1e^{s_1t}+A_2e^{s_2t}$, com $A_2=-A_1$ e
$A_1(s_1-s_2)=1{,}5$, resultando na expressão indicada.}
\end{questao}

\gabarito[Nível 4 --- fontes dependentes]
